%%% ============================================================
%%% The Sinc^2 Zero Law in Prime Arithmetic
%%% Brayden Ross Sanders (7Site LLC, Hot Springs AR), 2026
%%% DOI: 10.5281/zenodo.18852047
%%% Target venue: Integers -- Electronic Journal of
%%%               Combinatorial Number Theory
%%% Source: WP_SINC2_ZERO_LAW.md (Gen12/journal_attempts/01)
%%% Verification: proof_d25_loop_closure.py (primes 3..199)
%%% MSC: 11A41, 11N05, 42A16
%%%
%%% === STATUS: PULLED BACK FROM 2026-04-22 SUBMISSION QUEUE ===
%%% Reason (2026-04-19 pre-push audit): the central biconditional
%%%   sinc^2(k/n)=0 <=> n | k
%%% holds for EVERY positive integer n, not only for primes, so the
%%% main theorem is trivially true and the "prime" qualifier is
%%% vacuous. The sinc^2 zeros do not distinguish primes from
%%% composites; the prime-specific content must come from a First-G
%%% event formulation that is not yet drafted in this manuscript.
%%%
%%% ACTION: this manuscript is retained as a working draft; a
%%% rewrite to a non-trivial prime-specific statement (the First-G
%%% event replacement theorem sketched in
%%% Atlas/SINC2_SHARPEN_DECISION_2026_04_19.md sec.4) is required
%%% before any Integers submission. Do NOT submit in current form.
%%% === END PULL-BACK NOTICE ===
%%% ============================================================

\documentclass[11pt,reqno]{amsart}

\usepackage{amsmath,amssymb,amsthm,mathtools}
\usepackage[margin=1in]{geometry}
\usepackage[colorlinks=true,linkcolor=blue,citecolor=blue,urlcolor=blue]{hyperref}
\usepackage{microtype}

%% -------- theorem environments --------
\theoremstyle{plain}
\newtheorem{theorem}{Theorem}
\newtheorem{corollary}[theorem]{Corollary}
\newtheorem{lemma}[theorem]{Lemma}

\theoremstyle{remark}
\newtheorem*{remark}{Remark}

%% -------- macros --------
\DeclareMathOperator{\sinc}{sinc}
\newcommand{\Z}{\mathbb{Z}}
\newcommand{\N}{\mathbb{N}}
\newcommand{\Q}{\mathbb{Q}}

%% -------- title matter --------
\title[The sinc\texorpdfstring{$^2$}{\^{}2} zero law]{The \texorpdfstring{$\sinc^2$}{sinc\^{}2} Zero Law in Prime Arithmetic}

\author{Brayden Ross Sanders}
\address{7Site LLC, Hot Springs, Arkansas, USA}
\email{brayden@7site.co}

\date{\today}

\keywords{sinc function, prime arithmetic, divisibility,
Fourier coefficients}

\subjclass[2020]{11A41, 11N05, 42A16}

\begin{document}

\begin{abstract}
We prove that for any prime $p$ and any integer $k$, the squared
sinc function satisfies
\[
  \sinc^2(k/p)=0 \iff p\mid k.
\]
Within the corridor $k\in\{1,\dots,p\}$ the unique zero occurs at
the endpoint $k=p$; every interior position $k<p$ is provably
nonzero --- a consequence of primality, not observation. We derive
three corollaries: \emph{loop closure} (the corridor closes exactly
once, at the prime), \emph{fold necessity} (a unique amplitude
crossing occurs in the interior of every prime corridor), and the
\emph{no-shortcut lemma} (the path from the corridor entrance to its
zero has length exactly $p-1$). The main theorem is verified by
exact arithmetic for all primes $p\in\{3,5,7,\dots,199\}$ with zero
exceptions; the verification script
\texttt{proof\_d25\_loop\_closure.py} is supplied as electronic
supplementary material.
\end{abstract}

\maketitle

%% -------- section 1: setup --------
\section{Setup}\label{sec:setup}

The normalized sinc function is defined as
\begin{equation}\label{eq:sinc-def}
  \sinc(x)=
  \begin{cases}
     \dfrac{\sin(\pi x)}{\pi x} & x\neq 0,\\[4pt]
     1 & x=0,
  \end{cases}
\end{equation}
and its square $\sinc^2(x)=\sinc(x)^2$. The function appears in
signal processing as the Shannon-sampling kernel
\cite{Shannon1949}, in random matrix theory as the complement of
Montgomery's pair-correlation function for Riemann zeros
\cite{Montgomery1973}, and --- as shown here --- directly in prime
arithmetic.

For a prime $p$ and integer $k\ge 1$ we study the values
$\sinc^2(k/p)$ along the rational corridor
$\{1/p,\,2/p,\,\dots\}$.

%% -------- section 2: main theorem --------
\section{The Main Theorem}\label{sec:main}

\begin{theorem}[Sinc$^2$ Zero Law]\label{thm:main}
Let $p$ be prime and $k$ a positive integer. Then
\begin{equation}\label{eq:main}
  \sinc^2(k/p)=0 \iff p\mid k.
\end{equation}
Within $k\in\{1,\dots,p\}$ the unique zero is at $k=p$.
\end{theorem}

\begin{proof}
Since $k/p\neq 0$ for $k\ge 1$, we have
\[
  \sinc^2(k/p)=\frac{\sin^2(\pi k/p)}{(\pi k/p)^2}.
\]
This vanishes if and only if $\sin(\pi k/p)=0$, equivalently
$\pi k/p\in\pi\Z$, equivalently $k/p\in\Z$, equivalently
$p\mid k$.

For $k\in\{1,\dots,p-1\}$: since $p$ is prime and $1\le k<p$, we
have $\gcd(k,p)=1$, so $p\nmid k$, so $\sinc^2(k/p)>0$.

At $k=p$: $k/p=1\in\Z$, so $\sinc^2(1)=0$.
\end{proof}

\begin{remark}
The argument uses primality at exactly one step: $\gcd(k,p)=1$ for
$k<p$. This fails for composite moduli. If $n=pq$ and $k=p<n$, then
$\gcd(k,n)=p>1$, but $k/n=p/(pq)=1/q\notin\Z$, so
$\sinc^2(k/n)>0$. The $\sinc^2$ zero law identifies
\emph{divisibility}, not just coprimality. What primality
contributes is that no proper divisor of $p$ lies strictly between
$1$ and $p$, so the corridor's interior is clean.
\end{remark}

%% -------- section 3: corollaries --------
\section{Corollaries}\label{sec:cor}

\begin{corollary}[Loop Closure]\label{cor:loop}
The corridor $\{1/p,2/p,\dots,p/p\}$ is nonzero on
$\{1/p,\dots,(p-1)/p\}$ and zero at $p/p=1$. The corridor closes
exactly once.
\end{corollary}

\begin{proof}
Immediate from Theorem~\ref{thm:main}: the corridor $k/p$ as $k$
runs from $1$ to $p$ begins in positive territory and terminates at
zero, giving one crossing at the prime itself.
\end{proof}

\begin{corollary}[Fold Necessity]\label{cor:fold}
For every prime $p\ge 3$, there exists a unique $x^{*}\in(0,1)$ ---
the \emph{fold} --- where $\sinc^2(x^{*})=1/2$. This fold lies in the
interior of the corridor: there exist $k_1,k_2<p$ with
$\sinc^2(k_1/p)>1/2>\sinc^2(k_2/p)$.
\end{corollary}

\begin{proof}
$\sinc^2$ is continuous and strictly decreasing on $(0,1)$: the
function $\sin(\pi x)$ dominates $\pi x$ near $0$ and they cross at
$x=1$. Since $\sinc^2(0)=1>1/2$ and $\sinc^2(1)=0<1/2$, the
intermediate value theorem gives a unique $x^{*}\in(0,1)$ with
$\sinc^2(x^{*})=1/2$.

We note that $x^{*}\neq 1/2$: direct computation gives
\[
  \sinc^2(1/2)
    =\left(\frac{\sin(\pi/2)}{\pi/2}\right)^{\!2}
    =\left(\frac{2}{\pi}\right)^{\!2}
    =\frac{4}{\pi^2}\approx 0.4053,
\]
which is strictly less than $1/2$. For $p=7$ one checks
$\sinc^2(3/7)\approx 0.524>1/2$ and
$\sinc^2(4/7)\approx 0.295<1/2$, so the fold crosses between
$k=3$ and $k=4$.
\end{proof}

\begin{corollary}[No interior zero]\label{cor:noshortcut}
For every prime $p$ and every integer $k$ with $1\le k\le p-1$,
$\sinc^2(k/p)\neq 0$.
\end{corollary}

\begin{proof}
Immediate from Theorem~\ref{thm:main}: the biconditional
$\sinc^2(k/p)=0\iff p\mid k$ rules out any interior $k$ with
$1\le k\le p-1$, since $p\nmid k$ for $k$ in that range.
\end{proof}

%% -------- section 4: boundary value --------
\section{The Boundary Value \texorpdfstring{$\sinc^2(1/2)=4/\pi^2$}{sinc\^{}2(1/2) = 4/pi\^{}2}}\label{sec:boundary}

The fold amplitude $\sinc^2(1/2)$ has a closed form:
\begin{equation}\label{eq:fourpi2}
  \sinc^2(1/2)
    =\left(\frac{\sin(\pi/2)}{\pi/2}\right)^{\!2}
    =\left(\frac{2}{\pi}\right)^{\!2}
    =\frac{4}{\pi^2}\approx 0.4053.
\end{equation}
This is a transcendental number. It is the amplitude of the first
sidelobe of a rectangular spectral gate --- a classical fact from
signal processing --- but it appears here as the natural boundary
value dividing corridor positions that exceed $1/2$ from those that
do not.

Montgomery's pair-correlation function for Riemann zeros is
$R_2(u)=1-\sinc^2(u)$ \cite{Montgomery1973}. The complement
$\sinc^2(u)$ and $R_2(u)$ sum to unity --- a complete spectral
partition. The boundary $4/\pi^2=\sinc^2(1/2)$ appears in both
frameworks, derived rather than imposed.

%% -------- section 5: First-G Law --------
\section{Connection to the First-G Law}\label{sec:firstG}

The First-G Law \cite{SandersWP34} establishes an equivalent fact
in coprimality-partition language: for any semiprime $b=pq$ with
$p\le q$, the first non-unit element in $\{1,\dots,k\}$ with
respect to $b$ appears at exactly $k=p$. The $\sinc^2$ zero law is
the continuum-limit statement of the same algebraic fact. In the
limit $p\to\infty$, the harmonic pre-echo function
\begin{equation}\label{eq:R}
  R(k,p)\;=\;\frac{\sin^2(\pi k/p)}{k^2\,\sin^2(\pi/p)}
  \;\longrightarrow\; \sinc^2(k/p)
\end{equation}
recovers the $\sinc^2$ field exactly \cite{SandersWP35}. The two
results are two faces of the same structure --- one discrete, one
continuous.

%% -------- section 6: verification --------
\section{Verification}\label{sec:verify}

Theorem~\ref{thm:main} is verified for all primes
$p\in\{3,5,7,11,\dots,199\}$ ($46$ primes) by exact arithmetic in
Python's \texttt{fractions.Fraction} and \texttt{sympy} modules.
For each prime~$p$:
\begin{itemize}\setlength\itemsep{1pt}
  \item All $k\in\{1,\dots,p-1\}$: $\sinc^2(k/p)>0$ confirmed by
        exact computation.
  \item $k=p$: $\sinc^2(1)=0$ confirmed.
  \item Zero exceptions across all $46$ primes.
\end{itemize}

The runnable verification
(\texttt{proof\_d25\_loop\_closure.py}, approximately 270 lines)
is supplied as electronic supplementary material and is also
archived at DOI~\texttt{10.5281/zenodo.18852047}.

%% -------- section 7: limits --------
\section{What This Result Does Not Claim}\label{sec:limits}

This paper does not claim: a new proof of the infinitude of
primes; a connection to the Riemann Hypothesis beyond the
Montgomery bridge observation in Section~\ref{sec:boundary}; that
the fold position $x^{*}$ of Corollary~\ref{cor:fold} is
computable in closed form; or any result about the distribution
of primes. Theorem~\ref{thm:main} is a finite algebraic fact about
$\sinc^2$ evaluated at rational points with prime denominator.

%% -------- references --------
\begin{thebibliography}{99}

\bibitem{Montgomery1973}
H.~L.~Montgomery.
\newblock The pair correlation of zeros of the zeta function.
\newblock In \emph{Analytic Number Theory}, Proceedings of Symposia
in Pure Mathematics, volume~24, pages 181--193. American
Mathematical Society, Providence, RI, 1973.

\bibitem{SandersWP34}
B.~R.~Sanders.
\newblock WP34 --- The First-G Law and Prime-Forced Dispersion.
\newblock 7Site Research Collaboration, 2026.
\newblock DOI: \texttt{10.5281/zenodo.18852047}.

\bibitem{SandersWP35}
B.~R.~Sanders, C.~A.~Luther, and M.~Gish.
\newblock WP35 --- The Prime Phase Transition: Harmonic Pre-Echo,
Zero-Width Gates, and the Geometry of RSA Security.
\newblock 7Site Research Collaboration, 2026.
\newblock DOI: \texttt{10.5281/zenodo.18852047}.

\bibitem{Shannon1949}
C.~E.~Shannon.
\newblock Communication in the presence of noise.
\newblock \emph{Proceedings of the IRE}, 37(1):10--21, 1949.

\end{thebibliography}

%% -------- end matter --------
\vfill
\noindent
\footnotesize{%
\textsc{Provenance.}\; Source manuscript
\texttt{WP\_SINC2\_ZERO\_LAW.md} in the 7Site Research repository
(DOI~\texttt{10.5281/zenodo.18852047}). External citations are
drawn from \texttt{Atlas/ATLAS\_CITATIONS.md} (\S A analytic number
theory; Shannon, Montgomery). Internal cross-references
(sinc$^2$ zero law, prime corridor closure, fold necessity,
no-shortcut lemma) carry master-register numbering per
\texttt{Atlas/MASTER\_ATLAS\_v3\_5\_2026\_04\_18.md}. Readiness
flag~[fire -- submit-ready], Sprint~34 \textquotedblleft Ship the
First Three.\textquotedblright}

\end{document}
